\begingroup
\setlength{\LTcapwidth}{\textwidth}
\begin{longtable}{p{0.30\textwidth}p{0.26\textwidth}p{0.30\textwidth}}
\caption{Skenario verifikasi fungsional sistem.}
\label{tab:tab5.11_verifikasi_fungsional_sistem} \\
\toprule
\textbf{Skenario} & \textbf{Komponen} & \textbf{Keluaran yang diharapkan} \\
\midrule
\endfirsthead

\multicolumn{3}{@{}l}{\small\emph{Tabel \thetable{} (lanjutan)}} \\
\toprule
\textbf{Skenario} & \textbf{Komponen} & \textbf{Keluaran yang diharapkan} \\
\midrule
\endhead

\bottomrule
\endlastfoot

Memuat pasien & Frontend + API + penyimpanan & Data pasien dan logbook dapat ditampilkan \\

Membentuk prediksi & Prediction service & Prediksi sesuai modalitas dan horizon yang tersedia \\

Membentuk clinical state & Prediction + state & State terstruktur diteruskan ke pipeline berikutnya \\

Retrieval evidence & RAG retriever & Evidence dan metadata sumber dikembalikan \\

Generation assessment & LLM + evidence & Assessment terstruktur dihasilkan berdasarkan konteks \\

Citation resolution & Citation handler & Citation dapat ditelusuri ke evidence/sumber \\

Data tidak layak & Prediction service & Prediksi ditolak dan keterbatasan ditampilkan \\

Evidence tidak memadai & RAG pipeline & Keluaran menyatakan keterbatasan evidence \\

Integrasi CDSS & Frontend + backend & Hasil prediction dan assessment tampil pada alur yang sama \\
\end{longtable}
\endgroup
